\subsection{პროექტის თემასთან დაკავშირებული სფეროს კვლევა}

\subsubsection{Collaborative editing და არსებული ტექნოლოგიები}

რეალურ დროში ერთობლივი რედაქტირება დღეს ყველასთვის ნაცნობი ფუნქციაა -
Google Docs-მა ის მასობრივ პროდუქტად აქცია, მაგრამ ეს არ არის ერთადერთი ტექნოლოგია ამ სფეროში. 
თუ ამ ინსტრუმენტებს დავაკვირდებით, თითქმის ყველა მათგანი ერთსა და იმავე
არქიტექტურულ ნიმუშს მიჰყვება: არსებობს ცენტრალური სერვერი, რომელიც
დოკუმენტის „კანონიკურ“ ვერსიას ფლობს, ხოლო კლიენტები ამ სერვერს
უკავშირდებიან. ეს მიდგომა კარგად მუშაობს კომერციული SaaS პროდუქტისთვის,
მაგრამ მასთან ერთად მოდის დამოკიდებულებები: მუდმივი ინტერნეტი, მესამე
მხარის ინფრასტრუქტურა და დოკუმენტის მონაცემების cloud storage-ში შენახვა.

ტექნიკური თვალსაზრისით, პარალელური რედაქტირებისას conflict-resolution-ის ორი
ტექნიკა არსებობს. პირველი არის
\textbf{Operational Transformation (OT)} \cite{ot}: როდესაც ორი
მომხმარებელი ერთდროულად ასწორებს ტექსტს, მათი ოპერაციები (ჩასმა/წაშლა
პოზიციებით) ერთმანეთის მიმართ „ტრანსფორმირდება“ ისე, რომ განსხვავებული
თანმიმდევრობით შესრულების შემდეგაც ერთი და იგივე შედეგი მივიღოთ. OT
გამოიყენება Google Docs-სა და Etherpad-ში, თუმცა მას ორი ცნობილი სისუსტე
აქვს: ტრანსფორმაციის ფუნქციების კორექტულად დაწერა ძალიან რთულია 
და პრაქტიკული OT სისტემები თითქმის ყოველთვის ცენტრალურ სერვერს
საჭიროებს, რომელიც ოპერაციებს ერთიან თანმიმდევრობას მიანიჭებს.

მეორე მიდგომაა \textbf{CRDT, Conflict-free Replicated Data Type}
\cite{crdt}. აქ იდეა შებრუნებულია: იმის ნაცვლად, რომ ოპერაციები
ერთმანეთს მოვარგოთ, თავად მონაცემთა სტრუქტურა ისეა აგებული, რომ
ოპერაციები კომუტაციური იყოს, ანუ მიღების თანმიმდევრობას შედეგი არ
შეეცვალოს. თითოეული ჩასმული სიმბოლო (ან სიმბოლოთა ბლოკი) იღებს გლობალურად
უნიკალურ იდენტიფიკატორს და პოზიციის ნაცვლად სწორედ ამ იდენტიფიკატორებით
აღიწერება „სად“ ჩაჯდა ტექსტი. ამის ფასი ისაა, რომ წაშლილი ელემენტების
კვალი, ე.წ. tombstone-ები, გარკვეული ხნით უნდა შევინახოთ, სამაგიეროდ
ცენტრალური კოორდინატორი საერთოდ აღარ არის საჭირო: CRDT ბუნებრივად ერგება
peer-to-peer და დეცენტრალიზებულ არქიტექტურებს.

ტექსტისთვის განკუთვნილი არაერთი CRDT ალგორითმი არსებობს (WOOT, RGA,
Logoot/LSEQ და სხვა), მაგრამ პრაქტიკაში ყველაზე წარმატებული აღმოჩნდა
\textbf{YATA, Yet Another Transformation Approach} \cite{yata}, რომელიც
Yjs ბიბლიოთეკის \cite{yjs} საფუძველია. Yjs დღეს ალბათ ყველაზე ფართოდ
გამოყენებადი open source ერთობლივი რედაქტირების ძრავაა და სწორედ მისმა
პრაქტიკულმა წარმატებამ დაგვარწმუნა, რომ YATA სწორი არჩევანია. YATA-ში
ყოველი ჩასმა იმახსოვრებს არა მხოლოდ საკუთარ იდენტიფიკატორს, არამედ
„წარმოშობის“ მეზობლებსაც (origin left/right), თუ რომელ ორ ელემენტს შორის
დაიბადა ის. როდესაც ორი მონაწილე ერთსა და იმავე ადგილას ერთდროულად ჩასვამს
ტექსტს, კონფლიქტი დეტერმინისტული წესით წყდება და, რაც მნიშვნელოვანია,
პარალელურად აკრეფილი სიტყვები ერთმანეთში არ „აირევა“.

\subsubsection{რატომ CRDT და რატომ YATA}

ჩვენი არჩევანი სამ არგუმენტს ეყრდნობოდა. ჯერ ერთი, პროექტის მთავარი
მოტივაცია სწორედ ცენტრალურ სერვისზე დამოკიდებულების მოხსნა იყო, OT-ის
სერვერცენტრული ბუნება ამ მიზანს პირდაპირ ეწინააღმდეგება, CRDT კი
საშუალებას გვაძლევს, სერვერული კომპონენტი უბრალო relay-დ
ვაქციოთ, რომელსაც დოკუმენტის შინაარსი საერთოდ არ ესმის. მეორე, CRDT-ის
კორექტულობის მსჯელობა ლოკალურია: საკმარისია დავამტკიცოთ, რომ ოპერაციების
ინტეგრაცია კომუტაციურია, და ქსელის ნებისმიერი გაჭიანურება თუ
თანმიმდევრობის დარღვევა უსაფრთხო ხდება. მესამე, YATA-ს აქვს რეალურ
პროდუქტში (Yjs) წლების განმავლობაში გამოცდილი დიზაინი, საიდანაც ბევრი
პრაქტიკული ხერხი გადმოვიღეთ: სიმბოლოების ნაცვლად ბლოკებით (რამდენიმე
სიმბოლოს გაბმული მონაკვეთებით) ოპერირება და ბლოკის გაყოფა (splitting)
შუაში ჩასმისას, წაშლების კომპაქტური, დიაპაზონებად შენახვა და ა.შ.
დამატებით შთაგონებას ვიღებდით diamond-types პროექტიდანაც \cite{diamond},
რომელმაც გვიჩვენა, რომ CRDT დოკუმენტს პოზიციური ინდექსით მნიშვნელოვნად
აჩქარება შეუძლია.

\subsubsection{რატომ Cloudflare Tunnel}

დეცენტრალიზებული მიდგომის ერთი მოუხერხებელი დეტალი ისაა, რომ ჰოსტის
კომპიუტერი ინტერნეტიდან პირდაპირ მიუწვდომელია: სახლისა თუ უნივერსიტეტის
ქსელებში NAT და firewall გარე შეერთებებს ბლოკავს, ხოლო პორტების ხელით
გახსნა (port forwarding) რიგითი მომხმარებლისთვის შეუძლებელი მოთხოვნაა.
ამ პრობლემას ვჭრით Cloudflare Tunnel-ით \cite{cloudflared}: ჰოსტის
აპლიკაცია \texttt{cloudflared} პროგრამას უშვებს, რომელიც ჰოსტიდან
Cloudflare-ის ქსელისკენ \emph{გამავალ} კავშირს ამყარებს და სანაცვლოდ
დროებით საჯარო მისამართს (\texttt{*.trycloudflare.com}) იღებს. გამავალი
კავშირი NAT-ს უპრობლემოდ გადის, მომხმარებელს არაფრის კონფიგურაცია არ
სჭირდება, კავშირი TLS-ითაა დაცული და სერვისი უფასოა. თუ გვირაბის გახსნა
ვერ მოხერხდა, აპლიკაცია ამაზე არ ვარდება, სესია ლოკალური ქსელის ბმულით
გრძელდება.

\subsubsection{პროექტის ბუნება, მოთხოვნები და კომერციალიზაცია}

პროექტი არსებულ სფეროში არსებული გამოწვევის, ცენტრალიზებულ
ინფრასტრუქტურაზე დამოკიდებულების, ალტერნატიულ გზას გვთავაზობს.
ჩვენ არ გამოგვიგონია ახალი თანხმობის ალგორითმი; სიახლე კომბინაციაშია:
local-first desktop აპლიკაცია, რომელიც სესიის ინფრასტრუქტურას (relay
სერვერს და public tunnel-ს) მომხმარებლისვე კომპიუტერზე, ავტომატურად და
შეუმჩნევლად მართავს. ცალკე აღნიშვნის ღირსია ჩვენი garbage collection 
სისტემა, რომელიც, CRDT იმპლემენტაციებში გავრცელებული
მიდგომებისგან განსხვავებულ, Host-Authoritative პროტოკოლს იყენებს (დეტალურად -
ოპტიმიზაციების თავში).
